\chapter{Operators}\label{app-operators}

\section*{Every operator}

\begin{center}
\begin{longtable}{lP{5.4cm}P{4.6cm}}
\toprule
Operator & Means & Works on \\
\midrule
\endhead
\verb|a + b| & add; join text; join lists & numbers, quantities, \verb|str|, lists \\
\verb|a - b| & subtract & numbers, quantities, pointers \\
\verb|-a| & negate & numbers, quantities \\
\verb|a * b| & multiply & numbers; a quantity by a number \\
\verb|a / b| & divide, always giving a \verb|float| & numbers, quantities \\
\verb|a // b| & divide and round down & numbers \\
\verb|a % b| & remainder of \verb|//|, with the sign of \verb|b| & numbers \\
\verb|a ** b| & power & numbers \\
\verb|a & b| & bit and & whole numbers \\
\verb!a | b! & bit or & whole numbers \\
\verb|a ^ b| & bit exclusive or & whole numbers \\
\verb|~a| & every bit turned over & whole numbers \\
\verb|a << n|, \verb|a >> n| & shift left, shift right & whole numbers \\
\verb|a == b|, \verb|a != b| & equal, not equal, by contents & values of one type \\
\verb|a < b|, \verb|<=|, \verb|>|, \verb|>=| & order; they chain: \verb|0 <= i < n| & numbers, quantities, \verb|str| \\
\verb|x in xs|, \verb|x not in xs| & is it there? & lists, text, dict keys \\
\verb|x is none|, \verb|x is not none| & is it nothing? & \verb|T or none| \\
\verb|a and b|, \verb|a or b|, \verb|not a| & logic, asking the right side only when needed & \verb|bool| \\
\verb|xs[i]| & item \verb|i|, from 0 & lists, text, dicts, pointers \\
\verb|x.name| & a field, a method, a module's name, a unit change & \\
\verb|f(a, b)| & a call & functions, classes, enum cases \\
\bottomrule
\end{longtable}
\end{center}

\section*{Which goes first}

From the loosest to the tightest:

\begin{center}
\begin{tabular}{ll}
\toprule
Operators & \\
\midrule
\verb|or| & loosest \\
\verb|and| & \\
\verb|not| & \\
\verb|==  !=  <  <=  >  >=  in  not in  is  is not| & comparisons, which chain \\
\verb!|! & \\
\verb|^| & \\
\verb|&| & \\
\verb|<<  >>| & \\
\verb|+  -| & \\
\verb|*  /  //  %| & \\
\verb|-x  ~x| & \\
\verb|**| & tightest; and it groups to the right \\
\bottomrule
\end{tabular}
\end{center}

The bit operators bind tighter than the comparisons, as in Python and
unlike C. So this test means what it looks as if it means,
\verb|(flags & MASK) == 0|:

\begin{zenofrag}
let ready = flags & MASK == 0
\end{zenofrag}

\section*{The = forms}

Every operator on two numbers has an \verb|=| form, which updates a
\verb|var|, a field or an item:

\begin{center}
\verb|+=  -=  *=  /=  //=  %=  **=  &=  |\verb!|=!\verb|  ^=  <<=  >>=|
\end{center}

\verb|x += y| means \verb|x = x + y|, with \verb|x| read first. On a
sized number the result is checked as it is stored, so
\verb|count += 1| on a \verb|u8| of 255 is an \verb|OverflowError|.
\verb|/=| is refused on a whole number, because \verb|/| gives a
\verb|float|. A list takes \verb|+=| to add another list to its end.

\section*{What raises}

\begin{center}
\begin{tabular}{ll}
\toprule
When & Raises \\
\midrule
a whole-number answer does not fit in 64 bits & \verb|OverflowError| \\
\verb|//| or \verb|%| by zero, or \verb|/| by zero & \verb|ZeroDivisionError| \\
a whole number to a negative power & \verb|ValueError| \\
a shift by less than 0 or more than 63 & \verb|ValueError| \\
a shift left that loses bits & \verb|OverflowError| \\
\verb|xs[i]| outside the list or text & \verb|IndexError| \\
\verb|d[key]| of a key that is not there & \verb|KeyError| \\
\bottomrule
\end{tabular}
\end{center}
